\documentclass[12pt]{article}
\usepackage[utf8]{inputenc}
\usepackage[T1]{fontenc}
\usepackage{hyperref}
\title{Zauzou Language Resources for K-12 Curriculum Development}
\author{Compiled by Claude Sonnet 5.5 (claude-sonnet-5-5) for Andr\'as Kornai}
\date{}
\begin{document}
\maketitle

\section{Introduction}
Zauzou (ISO 639-3: zal; also called Rouruo) is a Loloish Sino-Tibetan language spoken in Tu'e District, Lanping County, Yunnan, China. It is most closely related to Nusu. Estimates of the number of speakers are around 2,100 (Bradley 2007), and other sources give different figures. It is not known to be taught in schools. Little published material exists, and no school materials were found in the searches.

\section{Primary Educational Resources}
No primers, textbooks or school curricula for Zauzou were found in the searches.

\section{Linguistic and Reference Materials}
\begin{itemize}
\item Wikipedia overview of the Zauzou language, with references including the Sun (2002) word lists for the Guoli and Jiangmo dialects and a 2019 book by Ch\'en H\v{a}ih\'ong on Zauzou in Lanping, Yunnan: \url{https://en.wikipedia.org/wiki/Zauzou_language}
\item ELDP project ``Documentation of Zauzou, an endangered language in China'' (Yu Li, Jiangmo dialect). Its outputs include a reference grammar and a word list: \url{https://cultureincrisis.org/projects/documentation-of-zauzou-an-endangered-language-in-china}
\item The ELDP archive collection for the same project: \url{https://eldp.access.preservica.com/dk0482/}
\end{itemize}

\section{Religious and Community Resources}
No Zauzou scripture or community-organisation resources were found in the searches. A people-group profile exists: \url{https://people-groups.asiaharvest.org/China/chinaPeoples/Z/Zauzou.pdf}

\section{Digital and Audio Resources}
The ELDP archive collection above holds the documentation materials, but I could not check which recordings are openly accessible. The Endangered Languages Project page for Zauzou is a general catalogue entry: \url{https://www.endangeredlanguages.com/lang/5168}

\section{Curriculum Development Context}
The language is vulnerable and no known school use exists. The documentation corpus and the descriptive grammar would be the starting points for an orthography and a lexicon, but they would have to be adapted with the community. No ready-made pedagogical material was found.

\end{document}